\section*{Summary}\addcontentsline{toc}{section}{Summary}

\Cref{tab:9-atlas} pairs the 24 kinds of excitation that this book treats with six classes of dark-matter coupling: electron density, nuclei, spin and current, vector absorption, axion-like and scalar absorption, and inelastic and other kinematics. The columns come from the target operators of \cref{sec:8.1}, the rows from the solid-state chapters. Of the 144 cells, 57 hold a channel, each with one home section; eight more are allowed at leading order but computed in no section, and 79 are empty at leading order. Four reasons account for nearly all of the empty ones. The first is operator content: the nuclear density creates no electronic excitation in the Born--Oppenheimer approximation, and the electron density no magnon without spin--orbit coupling. The others are the kinematics of absorption at $q = m_\phi v \simeq 0$, which no acoustic phonon or free nucleus can meet, the order in the coupling, and the absence of a new response for inelastic kinematics. Each row sits in a sector of the decomposition theorem. The four electronic rows of the density column are parts of one zero-phonon response; phonons, multiphonons and free recoil are electron-following and differ only in their lattice regime; the phonon-assisted row and part of the Migdal row are phonon-dressed. The dark photon, which couples to every charge, populates more of the atlas than any other model, and its loss-function rate on electrons and its polar-phonon rate are two sectors of one charge response.

The lightest mass that reaches a channel, \cref{eq:9.2-mch}, depends on the largest laboratory speed of the halo, the threshold and the shape of the dispersion. No channel lies below the floor $2\omega_{\mathrm{th}}/v_{\mathrm{max}}^2$, which a flat onset reaches. A gapless linear branch of speed $c_{\mathrm{ch}}$ lies above the floor by about $v_{\mathrm{max}}/4c_{\mathrm{ch}}$, because the particle must reverse its direction to hand over the momentum of the excitation: the acoustic edge of silicon at a threshold of \SI{1}{meV} is \SI{6.12}{keV} (\cref{we:9-acoustic}). A branch faster than $v_{\mathrm{max}}$, such as the bulk plasmon of silicon, is closed to elastic scattering. An offset minimum puts the electron--hole pairs of silicon between the phonon-assisted onset just above the floor, \SI{0.333}{MeV}, and the zero-phonon edge at 0.405 to \SI{0.408}{MeV}. A free nucleus needs a mass that grows as the square root of the threshold; a free silicon nucleus takes up the gap from \SI{46.8}{MeV}. Absorption reaches $m_\phi = \omega$, five and a half decades below scattering at equal energy. \Cref{fig:9-kinmap} and \cref{tab:9-edges} place silicon, germanium, GaAs, aluminum and ZrTe$_5$ on one map, which ranks reach and not rate. The edges move with $v_{\mathrm{max}}$ as \cref{eq:9.2-halo} states: as $v_{\mathrm{max}}^{-2}$ for a flat onset, nearly as $v_{\mathrm{max}}^{-1}$ for an acoustic edge.

On the plane of momentum transfer and energy deposit, \cref{fig:9-qomega} draws the support of every excitation of silicon, colored by sector. A pair without a phonon must take the whole momentum transfer and lives above $E_{\mathrm{min}}(\vect{q})$; a pair with a phonon starts just above the gap at every $q$, and between the two lies the off-shell window. Overlap on the plane is not interference. Sectors add, and different phonon numbers add; amplitudes interfere only for a coupling to electrons and nuclei at once and at the boundary between the two electronic sectors (\cref{tab:9-overlaps}). The only place where the cross term decides a rate at leading order is the electron-following sector, where the cloud that follows a silicon nucleus cancels its charge at long wavelength and leaves a dark photon without a single-phonon channel.

A Limit Check compares a published formula, written in the book's notation, with a fuller result one step further along one axis of the truncation label, or in a parameter that the label records in words, against a controlling parameter, and shades where the two agree to a stated tolerance. The axis of the step fixes the leading correction and the controlling parameter: $x_R$ for an impulse formula, $\xi_D$ for a Migdal probability at dipole order, and the momentum or energy along which screening changes for a formula on the many-body axis. Of the 75 Limit Checks of the channel chapters, 25 have their leading correction on a component of the label (\cref{fig:9-axes}, \cref{tab:9-limitchecks}), most of them in the electronic, lattice and Migdal chapters. The other 50, which include every magnetic one, expand in a parameter of their target, kinematics or coupling that the label records in words. \Cref{app:F} lists each formula with its source equation and the size of its leading correction. Of its 90 rows for these chapters, 80 are tested in the 75 boxes, and six carry a correction that the book does not size.
